\clearpage
\begingroup
\fancyhead[L]{\small THEORIE DER BEWEGUNG DES MONDES.}
\renewcommand{\headrulewidth}{0pt}
\SourceLeaf{632}
\section*{VIERTER ABSCHNITT.}
\begin{center}
\itshape Zweite Annäherung zur Breite.
\end{center}

\par\noindent\hspace*{\parindent}Bei der Berechnung der \textls[120]{Breite} lassen wir in \(\omega\) und \(\pi\) die Theile, die
den Factor \(\dfrac{1}{Rp}\) enthalten, weg und substituiren in dem die Perturbationskräfte
enthaltenden Theile der dritten Gleichung statt \(\theta\) seinen oben gefundenen
Werth \(g\sin\lambda+\tfrac38ng(1+\tfrac14n)\sin(2E-\lambda)\) und für
\(\dfrac{\mathrm{d}\lambda}{\mathrm{d}v}\), \(1+\tfrac34nn\). Hiedurch wird die Gleichung
\[
\begin{aligned}
0={}&\frac{\mathrm{dd}\theta}{\mathrm{d}v^2}
+\theta\left(1+\frac32\frac{nn}{1+\mu}\frac{B^3}{R^3}
\frac{\mathrm{d}u^2}{\mathrm{d}v^2}\right)\\
&+\frac32\frac{nn}{1+\mu}g\frac{B^3}{R^3}\frac{\mathrm{d}u^2}{\mathrm{d}v^2}\Bigl\{\\
&\quad -(1+\tfrac38nn)\sin(2(v-V)-\lambda)
-\tfrac38nn\sin(2(v-V)+\lambda)\\
&\quad -\tfrac38n(1-\tfrac34n)\sin(\lambda+2nv-2V)\\
&\quad +\tfrac38nn\sin(4v-2V-2nv-\lambda)\Bigr\}.
\end{aligned}
\]

\par\noindent\hspace*{\parindent}Nun ist
\[
\begin{aligned}
V&=nu-2\varepsilon\sin a+\tfrac54\varepsilon\varepsilon\sin2a\\
 &=nv-2\varepsilon\sin a+\tfrac54\varepsilon\varepsilon\sin2a+2ne\sin M,
\end{aligned}
\]
\noindent wo die Glieder vom 3\textsuperscript{ten} Grade übergangen sind;
\(a\) bedeutet übrigens \(nu-\mathrm{Apog.}\,\odot\), die mittlere Anomalie der Sonne;
wir werden es aber hier als \(nv-\mathrm{Apog.}\,\odot\) betrachten, da der daraus
entstehende Unterschied nur von dem Grade derjenigen Glieder ist, die wir
hier ohnehin vernachlässigen. Substituirt man diesen Werth von \(V\), so wird
\[
\begin{aligned}
2(v-V)-\lambda={}&2E-\lambda+4\varepsilon\sin a
-\tfrac52\varepsilon\varepsilon\sin2a-4ne\sin M
\end{aligned}
\]
\[
\begin{aligned}
\sin(2(v-V)-\lambda)={}&\sin(2E-\lambda)\mathbin{.}
(1-4\varepsilon\varepsilon+4\varepsilon\varepsilon\cos2a)\\
&+\cos(2E-\lambda)\mathbin{.}
(4\varepsilon\sin a-\tfrac52\varepsilon\varepsilon\sin2a-4ne\sin M)
\end{aligned}
\]
\noindent und hienach unsere Gleichung
\SourceLeaf{633}
\[
\begin{aligned}
0={}&\frac{\mathrm{dd}\theta}{\mathrm{d}v^2}
+\theta\left(1+\frac32\frac{nn}{1+\mu}\frac{B^3}{R^3}
\frac{\mathrm{d}u^2}{\mathrm{d}v^2}\right)\\
&+\frac32\frac{nn}{1+\mu}g\frac{B^3}{R^3}\frac{\mathrm{d}u^2}{\mathrm{d}v^2}\Bigl\{\\
&\quad -(1+\tfrac38nn-4\varepsilon\varepsilon)\sin(2E-\lambda)
-\tfrac38nn\sin(2E+\lambda)\\
&\quad -\tfrac38n(1-\tfrac34n)\sin\lambda+\tfrac38nn\sin(4E-\lambda)\\
&\quad +2ne\sin(2E+M-\lambda)-2ne\sin(2E-M-\lambda)\\
&\quad -\tfrac34n\varepsilon\sin(\lambda+a)+\tfrac34n\varepsilon\sin(\lambda-a)\\
&\quad -2\varepsilon\sin(2E-\lambda+a)+2\varepsilon\sin(2E-\lambda-a)\\
&\quad -\tfrac34\varepsilon\varepsilon\sin(2E-\lambda+2a)
-\tfrac{13}{4}\varepsilon\varepsilon\sin(2E-\lambda-2a)\Bigr\}.
\end{aligned}
\]
\noindent Ferner ist
\[
\frac{B^3}{R^3}=1+\tfrac32\varepsilon\varepsilon-3\varepsilon\cos a
+\tfrac92\varepsilon\varepsilon\cos2a
\]
\noindent und
\[
\begin{aligned}
\frac{\mathrm{d}u^2}{\mathrm{d}v^2}={}&1+2ee+4e\cos M+5ee\cos2M
-\tfrac{11}{2}nn\cos2E\\
&+\tfrac{15}{2}ne\cos(2E-M)+gg\cos2\lambda;
\end{aligned}
\]
\noindent folglich
\[
\begin{aligned}
\frac{B^3}{R^3}\frac{\mathrm{d}u^2}{\mathrm{d}v^2}
={}&1+2ee+\tfrac32\varepsilon\varepsilon+4e\cos M+5ee\cos2M
-\tfrac{11}{2}nn\cos2E\\
&+\tfrac{15}{2}ne\cos(2E-M)+gg\cos2\lambda
-3\varepsilon\cos a+\tfrac92\varepsilon\varepsilon\cos2a\\
&-6e\varepsilon\cos(M+a)-6e\varepsilon\cos(M-a);
\end{aligned}
\]
\noindent also die Gleichung
\[
\begin{aligned}
0={}&\frac{\mathrm{dd}\theta}{\mathrm{d}v^2}
+\theta\left\{1+\frac32\frac{nn}{1+\mu}(1+2ee+\tfrac32\varepsilon\varepsilon)\right\}\\
&+\frac32\frac{nn}{1+\mu}g(1+2ee+\tfrac32\varepsilon\varepsilon)\times
\end{aligned}
\]
\begingroup\small
\setlength{\arraycolsep}{3pt}\renewcommand{\arraystretch}{1.5}
\[
\begin{array}{r|l@{\qquad}r|l}
\multicolumn{1}{c}{\sin}& &\multicolumn{1}{c}{\sin}&\\
2E-\lambda&-1+\tfrac{19}{8}nn+4\varepsilon\varepsilon
 &\lambda+a&-\tfrac32\varepsilon+\tfrac{3}{16}n\varepsilon\\
2E+\lambda&-\tfrac{25}{8}nn-\tfrac12gg
 &\lambda-a&-\tfrac32\varepsilon+\tfrac{21}{16}n\varepsilon\\
\lambda&-\tfrac38n-\tfrac{79}{32}nn-\tfrac12gg
 &\lambda+2a&+\tfrac94\varepsilon\varepsilon\\
4E-\lambda&+\tfrac{25}{8}nn
 &\lambda-2a&+\tfrac94\varepsilon\varepsilon\\
2E+M-\lambda&-2e+\tfrac{11}{4}ne
 &2E-\lambda+a&-\tfrac12\varepsilon-\tfrac{9}{16}n\varepsilon\\
2E-M-\lambda&-2e-5ne
 &2E-\lambda-a&+\tfrac72\varepsilon-\tfrac{9}{16}n\varepsilon\\
2E-M+\lambda&+\tfrac{15}{4}ne
 &2E-\lambda+2a&0
\end{array}
\]
\endgroup
\noindent\hfill79\textsuperscript{*}\par
\SourceLeaf{634}
\begingroup\small
\setlength{\arraycolsep}{3pt}\renewcommand{\arraystretch}{1.5}
\[
\begin{array}{r|l@{\qquad}r|l}
\multicolumn{1}{c}{\sin}& &\multicolumn{1}{c}{\sin}&\\
2E-\lambda-2a&-\tfrac{17}{2}\varepsilon\varepsilon
 &4E-M-\lambda&-\tfrac{15}{4}ne\\
M+\lambda&+2e-\tfrac34ne
 &2E-3\lambda&-\tfrac12gg\\
M-\lambda&-2e-3ne
 &2M+\lambda&+\tfrac52ee\\
2E-\lambda+M+a&-e\varepsilon
 &2M-\lambda&-\tfrac52ee\\
2E-\lambda+M-a&+7e\varepsilon
 &3\lambda&+\tfrac12gg\\
2E-\lambda-M+a&-e\varepsilon
 &M+\lambda+a&-3e\varepsilon\\
2E-\lambda-M-a&+7e\varepsilon
 &M+\lambda-a&-3e\varepsilon\\
2E+2M-\lambda&-\tfrac52ee
 &M-\lambda+a&+3e\varepsilon\\
2E-2M-\lambda&-\tfrac52ee
 &M-\lambda-a&+3e\varepsilon.
\end{array}
\]
\endgroup
\noindent Da nach Integration dieser Gleichung das Glied \(\sin\lambda\) den Coefficienten
\(g\) haben muss, so ergibt sich durch eine ähnliche Operation, wie oben beim
Apogäum
\[
\frac{\mathrm{d}\lambda}{\mathrm{d}v}
=\BareRoot{\left\{1+\tfrac32\frac{nn}{1+\mu}
(1+2ee+\tfrac32\varepsilon\varepsilon-\tfrac12gg-\tfrac38n-\tfrac{79}{32}nn)\right\}},
\]
\noindent folglich bis auf die Grössen der 4\textsuperscript{ten} Ordnung
\[
\left[\frac{\mathrm{d}\lambda}{\mathrm{d}v}\right]
=1+\tfrac34\frac{nn}{1+\mu}
(1-\tfrac38n-\tfrac{91}{32}nn+2ee+\tfrac32\varepsilon\varepsilon-\tfrac12gg).
\]
\clearpage
\endgroup
